\subsection{Why a group becomes the baseline}
\label{subsec:grpo-group-advantage}

Suppose a response receives reward $8$. That number alone does not say whether the model
should repeat the response. If the model usually scores $9$ on this prompt, the response
was worse than expected; if it usually scores $2$, the same score was unusually good.
Learning needs a relative signal: reward minus the expected reward for this particular
prompt. This difference is the advantage.

One way to estimate the expected reward is to train a value model. GRPO instead estimates
it by direct sampling. At the beginning of a rollout iteration, freeze the current model
as the behavior policy $\pi_{\mathrm{old}}$ and sample $G$ sibling responses to the same
prompt:
\begin{equation}
  o_1,\ldots,o_G \sim \pi_{\mathrm{old}}(\cdot\mid q).
  \label{eq:grpo-grouped-rollouts}
\end{equation}
A verifier or reward model assigns one scalar $r_i=R(q,o_i)$ to each complete response.
The group mean is a Monte Carlo estimate of what this policy normally earns on this
prompt. Subtracting it makes average performance neutral, better responses positive, and
worse responses negative.

GRPO also divides by the group's reward standard deviation. The motivation is scale:
reward differences of $0.1$ may be substantial for one prompt while another prompt spans
the full interval from $0$ to $1$. Standardization makes the signal describe relative
standing within the group rather than inherit the raw scale of each prompt. This division
is a normalization heuristic, not a requirement of an unbiased baseline, and it changes
the weighting induced by different groups.

These choices give the group-relative advantage
\begin{equation}
  A_i
  = \frac{r_i-\mu_q}{\sigma_q+\epsilon_A},
  \qquad
  \mu_q=\frac{1}{G}\sum_{j=1}^{G}r_j,
  \qquad
  \sigma_q^2=\frac{1}{G}\sum_{j=1}^{G}(r_j-\mu_q)^2.
  \label{eq:grpo-group-advantage}
\end{equation}

For example, suppose four responses to ``$2+3$'' receive rewards
$\{1.1,1.1,0.1,0.1\}$. Their mean is $0.6$ and their population standard deviation is
$0.5$. The two correct responses receive advantages near $+1$; the two incorrect
responses receive advantages near $-1$. The group acts like a small tournament among
answers to the same question.

This replacement has a clear tradeoff. It removes the need to train a separate critic for
the baseline, but it pays for several generations per prompt and retains Monte Carlo
sampling noise. If all siblings receive the same reward, the group contains no evidence
about which behavior was better: the relative learning signal is zero, and implementations
must handle the zero-variance case explicitly. Finally, $A_i$ still belongs to the whole
response. It has not yet explained how one number changes thousands of token decisions;
that is the next problem.
